\section*{第 \ref{ch:b3:measure} 章——测度论}

\begin{solution}{exo:b3:measure:1}
(a) 取补运算交换了两种定义情形。可数集的
可数并仍为可数集；若其中某一成员为余可数集，
则其并为余可数集：稳定性成立。它
包含全部单点集，而任一包含它们的
$\sigma$-代数都包含全部可数集（即可数并）及其
补集：故它即为 $\sigma(\{\text{单点集}\})$。

(b) 记 $\mathcal B = \sigma(\text{开集})$。
$\R$ 的每个开子集都是有理端点开区间的可数并
（对\omterm{def:b3:topology:topology}{开集}中每个有理点，取一个半径为有理数
且含于该\omterm{def:b3:topology:topology}{开集}的区间），故\omterm{def:b3:topology:topology}{开集} $\in
\sigma(\text{开区间}) \subseteq \sigma(\text{有理
端点区间})$。互相表示如下：$\intoo ab = \bigcup_n\intcc{a +
\frac1n}{b - \frac1n}$；$\intcc ab = \bigcap_n \intoo{a -
\frac1n}{b + \frac1n}$；$\intoc{-\infty}a = \bigcap_n
\intoo{-\infty}{a + \frac1n}$，反之 $\intoo ab =
\intoo{-\infty}b \setminus \intoc{-\infty}a$；有理射线：
$\intoc{-\infty}a =
\bigcap_{q \in \Q,\, q > a}\intoc{-\infty}q$。这四族中每一族
经可数次运算即可生成其余各族：四者皆生成
$\mathcal B$。

(c) 若只允许有限端点，则此族甚至不是
集代数：$\intoc01$ 的补集含有无界射线。
若允许无穷端点（$\intoc{-\infty}b$、$\intoo
a{+\infty}$），则它成为一个集代数（补运算及有限并
运算仍保持此类集合），但不是 $\sigma$-代数：
$\{0\} = \bigcap_n\intoc{-\frac1n}0$ 不是有限个
非退化区间的并。
\end{solution}

\begin{solution}{exo:b3:measure:2}
(a) $A \cup B = A \sqcup (B \setminus (A\cap B))$，故
$\mu(A\cup B) = \mu(A) + \mu(B) - \mu(A\cap B)$（有限性
允许作减法）。三个集合时，对两集合公式
连用两次，
\[
\mu(A\cup B\cup C) = \sum\mu(A) - \sum\mu(A\cap B) +
\mu(A\cap B\cap C)
\]
（各求和号的指标集显然）。

(b) 对\omterm{def:b3:measure:lebesgueouter}{勒贝格测度}，$A_n = [n, +\infty) \downarrow
\varnothing$，但 $\lambda(A_n) = \infty \not\to 0$。

(c) 一点落在长为 $\varepsilon$ 的区间内，
$\lambda(\{x\}) = 0$；可数次可加性使得可数
集\omterm{def:b3:measure:measure}{测度}为零。于是 $\lambda(\Q \cap \intcc01) = 0$，由可加性，
$\lambda(\intcc01\setminus\Q) = 1$：无理数承载了
全部长度。
\end{solution}

\begin{solution}{exo:b3:measure:3}
形如 $\intoc{-\infty}t$ 的射线构成一个 $\pi$-系（两者
之交即为较小者），生成 $\mathcal B(\R)$
（见\cref{exo:b3:measure:1}）。集合 $P_k = \intoc{-\infty}k$
递增趋于 $\R$，且 $\mu(P_k) \leq 1 < \infty$：
由\cref{thm:b3:measure:uniqueness}可知，$\mu = \nu$ 在
$\mathcal B(\R)$ 上成立。因此分布函数决定了
\omterm{def:b3:measure:measure}{测度}。
\end{solution}

\begin{solution}{exo:b3:measure:4}
对每个 $N$，$\limsup A_n \subseteq \bigcup_{n \geq N}A_n$，
故 $\mu(\limsup A_n) \leq \sum_{n\geq N}\mu(A_n)$，其右端为
收敛级数的尾项；令 $N \to \infty$。应用：取
$A_n = \{x \in \intcc01 : \abs{x - r_n} \leq 4^{-n}\}$
（$r_n$ 为第 $n$ 个有理数），$\lambda(A_n) \leq 2\cdot4^{-n}$
可和，于是 $\lambda(\limsup A_n) = 0$，即几乎每个 $x$
只属于有限个 $A_n$。（然而每个 $x$ 都是
有理数的极限：要点在于\emph{速度}$4^{-n}$。）
\end{solution}

\begin{solution}{exo:b3:measure:5}
$K = \bigcap K_n$ 是有限个闭区间之并的交，
故为紧集。第 $n$ 步余下 $2^n$ 个区间
等长，长为 $\ell_n \leq 2^{-n}$（每步
二等分并缩短）；若某区间 $I \subseteq K$，则它
须含于某一第 $n$ 步区间之内（对每个 $n$），迫使 $\ell(I) = 0$：
\omterm{def:b3:topology:interior}{内部}为空。移去的\omterm{def:b3:measure:measure}{测度}为 $\sum_{n\geq1}2^{n-1}\cdot
4^{-n} = \frac12\sum_{n\geq1}2^{-n} = \frac12$，所移去的区间
皆为互不相交的开区间，故 $\lambda(K) = \frac12$。

变式：移去长为 $\varepsilon
4^{-n}$ 的中央区间，余下一无处稠密的紧集 $K^{(\varepsilon)}$，
其\omterm{def:b3:measure:measure}{测度}为 $1 - \frac\varepsilon2$。于是 $\bigcup_m
K^{(1/m)}$ 为\omterm{rem:b3:complete:meagre}{贫乏集}（可数个无处\omterm{def:b3:topology:interior}{稠密集}之并），
\omterm{def:b3:measure:measure}{测度} $\geq \sup_m(1 - \frac1{2m}) = 1$：一个
满\omterm{def:b3:measure:measure}{测度}的\omterm{rem:b3:complete:meagre}{贫乏集}——且其在 $\intoo01$ 中的余集为
稠密的 $G_\delta$ 集，\omterm{def:b3:measure:measure}{测度}为 $0$（\omterm{def:b3:topology:topology}{拓扑}上“大”，度量上为零）。
$K^{(\varepsilon)}$ 在 $\intoo01$ 中的余集为
开的、稠密的，\omterm{def:b3:measure:measure}{测度}为 $\frac\varepsilon2 < \varepsilon$。
\end{solution}

\begin{solution}{exo:b3:measure:6}
将 $\intoc01$ 切成 $2^n$ 个 $\intoc0{2^{-n}}$ 的平移：
$c = 2^n\,\mu(\intoc0{2^{-n}})$，故 $\mu(\intoc0{2^{-n}}) =
c\,2^{-n} = c\,\lambda(\intoc0{2^{-n}})$。由平移不变性
及可加性，$\mu = c\lambda$ 在每个区间
$\intoc ab$ 上成立，其中 $b - a$ 为二进有理数、$a$ 任意；一般的
$\intoc ab$ 是此类区间的递增并
（$b_k \uparrow b$ 为自 $a$ 起的二进步长），再由下连续性
可将等式推广至一般情形。区间 $\intoc ab$ 构成一个
$\pi$-系，生成 $\mathcal B(\R)$，且 $\intoc{-k}k
\uparrow \R$ 且\omterm{def:b3:measure:measure}{测度}有限（$\mu(\intoc{-k}k) = 2kc$）：
由\cref{thm:b3:measure:uniqueness}可知 $\mu = c\lambda$ 在
$\mathcal B(\R)$ 上成立。
\end{solution}

\begin{solution}{exo:b3:measure:7}
由正则性（\cref{thm:b3:measure:regularity}）取 $K
\subseteq A \subseteq U$，其中 $K$ 紧、$U$ 开，
并使 $\lambda(U\setminus K) < \varepsilon$（两次逼近
均在 $\varepsilon/2$ 之内，且 $\lambda(U \setminus K) =
\lambda(U\setminus A) + \lambda(A \setminus K)$）。将 $U$ 写成
可数个互不相交的开区间 $(I_n)$ 之并（
即该\omterm{def:b3:topology:topology}{开集}的\omterm{def:b3:topology:components}{连通分支}）。紧集 $K$ 被其中有限个
开区间覆盖：$K \subseteq B = I_1\cup\dots\cup I_N \subseteq
U$。于是 $A \setminus B \subseteq A\setminus K \subseteq
U\setminus K$，且 $B \setminus A \subseteq U \setminus A
\subseteq U\setminus K$：
$\lambda(A\,\triangle\,B) \leq 2\lambda(U\setminus K)$——
取 $\varepsilon/2$ 出发即可使结果小于 $\varepsilon$。
\end{solution}

\begin{solution}{exo:b3:measure:8}
将 $A$ 换成\omterm{def:b3:measure:measure}{测度}为正的 $A \cap [-M, M]$（某个 $M$
可行，由自下连续性），不妨设 $0 < \lambda(A) <
\infty$。取\omterm{def:b3:topology:topology}{开集} $U \supseteq A$，使 $\lambda(U) <
\frac43\lambda(A)$，并将 $U = \bigsqcup_nI_n$ 分解为
互不相交的开区间：$\lambda(A) = \sum_n\lambda(A\cap
I_n)$。若对每个 $n$ 都有 $\lambda(A\cap I_n) \leq
\frac34\ell(I_n)$，求和即得 $\lambda(A) \leq
\frac34\lambda(U) < \lambda(A)$，矛盾。故存在某个区间 $I$ 满足
$\lambda(A\cap I) > \frac34\ell(I)$。令 $B = A \cap I$，并设
$\abs t < \frac12\ell(I)$，则 $B$ 与 $B + t$ 均落在
区间 $I \cup (I + t)$ 中，其长度 $< \frac32\ell(I)$。若
二者不交，则 $\lambda(B) + \lambda(B + t) = 2\lambda(B)
> \frac32\ell(I)$，超出包含区间的\omterm{def:b3:measure:measure}{测度}，
矛盾。故 $B \cap (B + t) \neq \varnothing$：
即存在 $x \in B$ 写成 $x = y + t$，其中 $y \in B$，从而 $t = x -
y \in A - A$。
因此 $\intoo{-\frac{\ell(I)}2}{\frac{\ell(I)}2} \subseteq A -
A$。
\end{solution}

\begin{solution}{exo:b3:measure:9}
维塔利平移族 $(V + q)_{q\in\Q}$ 划分 $\R$（每个
实数恰与唯一的代表元等价）。假设所有
集合 $B_q = A \cap (V + q)$ 都可测。$V + q$ 中任意两元素之差
为无理数或零（两个
不同代表元不等价），故 $B_q - B_q$
与 $\Q$ 的交仅为 $\{0\}$：它不含任何区间，于是
由斯坦豪斯（\cref{exo:b3:measure:8}）推出 $\lambda(B_q) = 0$。
于是 $\lambda(A) \leq \sum_{q}\lambda(B_q) = 0$，与
$\lambda(A) > 0$ 矛盾。故某个 $B_q \subseteq A$ 不可测。
\end{solution}

\begin{solution}{exo:b3:measure:10}
\emph{可测 $\Rightarrow$ $\varepsilon$-逼近}：由
外正则性（\cref{thm:b3:measure:regularity}）取\omterm{def:b3:topology:topology}{开集}
$U \supseteq A$，使 $\lambda(U) \leq \lambda(A) +
\varepsilon$；由可测性可作减法
$\lambda(U\setminus A) = \lambda(U) - \lambda(A) \leq
\varepsilon$。\emph{$\varepsilon$-版本 $\Rightarrow$
$G_\delta$-版本}：取 $U_n$ 使 $\lambda^*(U_n\setminus
A) < \frac1n$，并令 $G = \bigcap U_n$，为 $G_\delta$ 集，且
$\lambda^*(G\setminus A) \leq \lambda^*(U_n\setminus A) \to 0$。
\emph{$G_\delta$-版本 $\Rightarrow$ 可测}：$G\setminus
A$ 为 $\lambda^*$-零测集，故由完备性
（\cref{thm:b3:measure:caratheodory}）可测，且 $A = G \setminus
(G\setminus A)$ 可测（$G$ 为博雷尔集）。因此勒贝格可测集恰为
“博雷尔集模零测集”。
\end{solution}

\begin{solution}{exo:b3:measure:11}
(a) $\liminf A_n = \bigcup_N\bigcap_{n\geq N}A_n$ 是
集合 $B_N = \bigcap_{n\geq N}A_n$ 的递增并，
故 $\mu(\liminf A_n) = \lim_N\mu(B_N)$（由自下连续性
）；又 $\mu(B_N) \leq \inf_{n \geq N}\mu(A_n)$，其极限为
$\liminf\mu(A_n)$。对 $\limsup$，对补集作同样论证，
在有限\omterm{def:b3:measure:measure}{测度}的母空间内进行：
$U = \bigcup A_n$——对
递减的 $C_N = \bigcup_{n\geq N}A_n$ 使用自上连续性，要求 $\mu(C_1)
\leq \mu(U) < \infty$，并得 $\mu(\limsup A_n) =
\lim\mu(C_N) \geq \limsup\mu(A_n)$。

(b) 移动区间 $A_n = \intcc n{n+1}$：每个点
至多属于两个 $A_n$，且最终不属于其中任何一个，
故 $\limsup A_n = \varnothing$；但 $\mu(A_n) = 1$。于是
$\limsup\mu(A_n) = 1 > 0 = \mu(\limsup A_n)$。若没有
有限\omterm{def:b3:measure:measure}{测度}的包络，(a) 中的第二个不等式便不成立
——质量逃逸到无穷远，那里没有任何固定集合
能将其捕获。

(c) 若 $\sum\mu(A_n) < \infty$，则 $\mu(C_N) \leq
\sum_{n\geq N}\mu(A_n) \to 0$，且 $\limsup A_n =
\bigcap C_N$ 的\omterm{def:b3:measure:measure}{测度} $\leq \inf\mu(C_N) = 0$。单调
情形：递增时为自下连续性；递减且
$\mu(A_1) < \infty$ 时为自上连续性——
两者均已在\cref{ch:b3:measure}的基本性质中证明；
反例 $A_n = \intco n\infty$（递减趋于
$\varnothing$，但 $\mu \equiv \infty$）表明有限性
同样是必不可少的。
\end{solution}

\begin{solution}{exo:b3:measure:12}
(a) 集合 $E_{k,N}$ 随 $N$ 递增（约束更少），
且每个 $x$ 最终满足 $\abs{f_n(x)
- f(x)} \leq \frac1k$ 对所有 $n \geq N(x)$ 成立（逐点
收敛）：$\bigcup_NE_{k,N} = X$。由自下连续性：
$\mu(E_{k,N}) \to \mu(X) < \infty$，故 $\mu(X\setminus
E_{k,N}) \to 0$；据此选取 $N_k$。

(b) 令 $A = \bigcap_kE_{k,N_k}$，则 $\mu(X\setminus A) \leq
\sum_k\varepsilon2^{-k} = \varepsilon$。在 $A$ 上：对每个
$k$，所有 $n \geq N_k$ 满足 $\sup_A\abs{f_n - f} \leq
\frac1k$——恰为 $A$ 上的一致收敛。

(c) 对移动凸起而言，$A$ 上的一致收敛要求
$A$ 最终避开每个 $\intcc n{n+1}$——更精确地说，
$\sup_A\abs{f_n} < \frac12$ 迫使 $A \cap \intcc
n{n+1}$ 对充分大的 $n$ 为空，因此 $X \setminus A$
包含尾部 $\bigcup_{n\geq n_0}\intcc n{n+1}$，其\omterm{def:b3:measure:measure}{测度}为无穷。
在 (a) 中，有限性将“$E_{k,N}\nearrow X$”
转化为“补集的\omterm{def:b3:measure:measure}{测度}
趋于 $0$”：自上连续性需要有限的起点，
而在无穷\omterm{def:b3:measure:measure}{测度空间}上，逃逸到无穷远
正是它所无法察觉的。
\end{solution}

\begin{solution}{pb:b3:measure:1}
\textbf{1.} 归纳法。连续性：三个公式在
接合点处取值相同（$\frac12c_n(1) = \frac12$ 与 $\frac12 +
\frac12c_n(0) = \frac12$）；每一段都连续。
单调性与边界值均被继承。关于
压缩估计：在 $[0,\frac13]$ 上，$\abs{c_{n+1} - c_n}(x)
= \frac12\abs{c_n - c_{n-1}}(3x) \leq \frac12\norm{c_n -
c_{n-1}}_\infty$；在中间三分之一段上差为 $0$；在
右三分之一段上，与左段相同。

\textbf{2.} $\norm{c_{n+1} - c_n}_\infty \leq
2^{-n}\norm{c_1 - c_0}_\infty$：增量构成的级数
一致收敛，故 $c_n \to c$ 一致收敛；
$c$ 连续、非减，$c(0) = 0$，$c(1) = 1$（这些性质
经一致极限保持），在定义递推式两边取极限，
可知 $c$ 本身满足
三个自相似恒等式。

\textbf{3.} 由中间恒等式，$c \equiv \frac12$ 在
$\intcc{\frac13}{\frac23}$ 即第一个缺口。$C$ 的每个缺口都是
第一个缺口经两个仿射压缩的复合之像，
即 $x \mapsto \frac x3$，$x\mapsto\frac{x +
2}3$；恒等式据此传递常值性（取值为二进有理数）。
关于数位公式，取 $x =
\sum_n 2b_n3^{-n} \in C$：若 $b_1 = 0$，则 $x \in
[0,\frac13]$ 且 $c(x) = \frac12c(3x)$，$3x$ 的数字为
$(b_2, b_3, \dots)$；若 $b_1 = 1$，则 $x \in [\frac23, 1]$
且 $c(x) = \frac12 + \frac12c(3x - 2)$，移位方式相同。
由归纳法，前 $N$ 位二进数字（$c(x)$ 的）依次为 $b_1,
\dots, b_N$，对每个 $N$ 均成立：$c(x) = \sum_nb_n2^{-n}$。

\textbf{4.} 在 $C$ 之外，$c$ 局部为常数，从而可导，
导数为 $0$。由 $\lambda(C) = 0$
（\cref{ex:b3:measure:cantor}）可知 $c' = 0$ 几乎处处成立。然而
$c(1) - c(0) = 1$：微积分基本定理的 $\mathcal
C^1$ 形式要求 $c$ \emph{处处}可导，且导数
连续（或至少可积并满足绝对连续性，见
\cref{ch:b3:lebesgue}）；$c$ 在 $C$ 的点处不可导，
且更根本地，$c$ 不满足
\emph{绝对连续性}：它在一个零测集上攀升。

\textbf{5.} 给定 $y = \sum_n\beta_n2^{-n} \in \intcc01$
（$\beta_n \in \{0,1\}$），点 $x = \sum_n2\beta_n3^{-n}
\in C$ 由第 3 题有 $c(x) = y$，故 $c(C) = \intcc01$，
一个\omterm{def:b3:measure:measure}{测度}为 $1$ 的集合——零测集 $C$ 经 $c$ 承载了
整个区间。

\textbf{6.} $h$ 连续，且严格递增（$x$ 严格递增，
$c$ 非减）；$h(0) = 0$，$h(1) = 1$，由介值定理，
$h$ 是 $\intcc01$ 上的连续双射；
\omterm{def:b3:topology:compact}{紧空间}到\omterm{def:b3:topology:hausdorff}{豪斯多夫空间}的连续双射，
必为\omterm{def:b3:topology:continuity}{同胚}
（见\cref{cor:b3:topology:compacthomeo}）。

\textbf{7.} 在缺口 $(u, v)$（长度为 $\ell$）上，$c$ 为常数，
故 $h$ 是斜率为 $\frac12$ 的仿射函数：$h((u,v))$ 为
长度为 $\frac\ell2$ 的区间。诸缺口两两不交，且 $h$
为单射：其像两两不交，总\omterm{def:b3:measure:measure}{测度}为
$\frac12\sum\ell = \frac12(1 - \lambda(C)) = \frac12$.

\textbf{8.} $h(\intcc01) = \intcc01$ 和 $h(C)$ 为紧集
（连续像），因而可测，并且
\[
\lambda\bigl(h(C)\bigr) = 1 -
\lambda\bigl(h(\intcc01\setminus C)\bigr) = 1 - \tfrac12 =
\tfrac12 .
\]
\omterm{def:b3:topology:continuity}{同胚}可以把零测集放大为\omterm{def:b3:measure:measure}{测度}为 $\frac12$ 的集合：
“\omterm{def:b3:topology:topology}{拓扑}大小”（范畴、维数）与“\omterm{def:b3:measure:measure}{测度}”
在\omterm{def:b3:topology:continuity}{同胚}下的传递方式截然不同，
只有前者是\omterm{def:b3:topology:topology}{拓扑}不变量。

\textbf{9.} $Z = h^{-1}(W) \subseteq h^{-1}(h(C)) = C$ 具有
$\lambda^*(Z) \leq \lambda(C) = 0$：是零测集，故
由完备性知其为勒贝格可测的
（见\cref{thm:b3:measure:caratheodory}）。

\textbf{10.} 设 $\varphi$ 连续，$\mathcal D = \{B
: \varphi^{-1}(B) \in \mathcal B\}$。原像运算与补集、可数并可交换，
故 $\mathcal D$ 是 $\sigma$-代数；由连续性，它包含\omterm{def:b3:topology:topology}{开集}，
因而 $\mathcal D \supseteq \mathcal B$——
即\omterm{def:b3:topology:continuity}{连续映射}下博雷尔集的原像
仍为博雷尔集。

\textbf{11.} 若 $Z$ 为博雷尔集，则对
\omterm{def:b3:topology:continuity}{连续映射} $\varphi = h^{-1}$ 应用第 10 题：$\varphi^{-1}(Z) = h(Z) = W$
将为博雷尔集，从而为勒贝格可测，与 $W$ 的选取矛盾。
故 $Z \in \mathcal L \setminus \mathcal B$：
勒贝格 $\sigma$-代数严格包含博雷尔集全体。

\textbf{12.} $g = \mathbf 1_Z$ 为勒贝格可测（$Z \in
\mathcal L$）和 $\varphi = h^{-1}$ 连续，但
$(g\circ\varphi)^{-1}(\{1\}) = \varphi^{-1}(Z) = W$ 不可测，
故 $g \circ \varphi$ 不是勒贝格可测的。这里须小心：
“勒贝格可测函数”是指：
\emph{博雷尔}集的原像落在 $\mathcal L$ 中；复合则要求
\emph{勒贝格}集的原像须仍为勒贝格集，这是
连续性所不能保证的（此处 $\varphi^{-1}(Z) \notin
\mathcal L$，尽管 $\varphi$ 是\omterm{def:b3:topology:continuity}{同胚}）。

\textbf{13.} 各链如下，并给出严格性的例子：
\[
\{\text{可数}\} \subsetneq \{\text{博雷尔零集}\}
\subsetneq \{\text{勒贝格零集}\} \subsetneq \mathcal L
\subsetneq \mathcal P(\R),
\qquad
\{\text{博雷尔零集}\} \subsetneq \mathcal B \subsetneq
\mathcal L .
\]
例：$C$ 为博雷尔集、零集，且不可数（对应第一处真包含）；$Z$ 为
勒贝格零集但不是博雷尔集（对应第二处真包含，以及在 $\mathcal B
\subsetneq \mathcal L$ 内的第六处）；一个胖康托集是博雷尔集，
无处稠密、\omterm{def:b3:measure:measure}{测度}为正（用于区分零集与
博雷尔集）；维塔利集 $V$ 不属于 $\mathcal L$（对应最后一处真包含）。
\omterm{def:b3:measure:measure}{测度}、\omterm{def:b3:topology:topology}{拓扑}与基数分别从
本质不同的角度切分 $\mathcal P(\R)$。

\textbf{14.} \omterm{def:b3:measure:outer}{外测度}：$\mu_F^*(\varnothing) = 0$（取一个
退化区间作覆盖即可），单调性显然，可数次可加性
由拼接各 $\varepsilon2^{-k}$ 最优覆盖得到，
与 $\lambda^*$ 的情形完全相同。单个区间的覆盖
给出 $\mu_F^*(\intoc ab) \leq F(b) - F(a)$。反之，设
$\intoc ab \subseteq \bigcup_k\intoc{a_k}{b_k}$。由
$F$ 的连续性，取 $b_k' > b_k$，使 $F(b_k') \leq F(b_k) +
\varepsilon2^{-k}$，并取 $a' \in \intoo ab$，使 $F(a') \leq
F(a) + \varepsilon$。紧集 $\intcc{a'}b$ 被
开区间 $\intoo{a_k}{b_k'}$ 覆盖：有限个即可，并
借助 \cref{thm:b3:measure:lebesgue} 中的链式论证（从
$a'$ 沿重叠区间走到 $b$，裂项
相消 $F$ 的增量，单调性吸收重叠部分），得 $F(b)
- F(a') \leq \sum_k(F(b_k') - F(a_k)) \leq \sum_k(F(b_k) -
F(a_k)) + \varepsilon$。令 $\varepsilon \to 0$，即得等式。
（$F$ 的连续性使得可以在
任意小的 $F$ 代价下打开区间。）

\textbf{15.} 只需证明每条半直线 $H_t =
\intoc{-\infty}t$ 是卡拉泰奥多里可测的，因为
可测集构成 $\sigma$-代数
（见\cref{thm:b3:measure:caratheodory}），且半直线生成
$\mathcal B$。设 $A$ 及其 $\varepsilon$ 最优覆盖
$(\intoc{a_k}{b_k})$ 为 $A$ 的覆盖，每个区间分解为
$\intoc{a_k}{t\wedge b_k} \cup \intoc{t \vee a_k}{b_k}$（其中一块
可能为空），两块的 $F$ 代价恰好相加为
$F(b_k) - F(a_k)$；前一块覆盖 $A \cap H_t$，后一块
覆盖 $A \setminus H_t$。因此 $\mu_F^*(A\cap H_t) +
\mu_F^*(A\setminus H_t) \leq \mu_F^*(A) + \varepsilon$，而
反向不等式即次可加性。将所得
\omterm{def:b3:measure:measure}{测度}限制到 $\mathcal B$ 上，即得勒贝格–斯蒂尔杰斯
\omterm{def:b3:measure:measure}{测度} $\mu_F$。

\textbf{16.} $\mu(\R) = \lim_n(F(n) - F(-n)) = 1 - 0 = 1$
（由\omterm{def:b3:measure:measure}{测度}沿 $\intoc{-n}n$ 的连续性）。
设 $\intoo uv$ 为 $C$ 的一个缺口，则 $c$ 为常数，故每个半开
子区间为 $\mu$ 零集，缺口本身亦然（由可数
并得到）；在 $\intcc01$ 之外 $c$ 同样为常数。因此
$\mu(\R\setminus C) = 0$，$\mu(C) = 1$，而 $\lambda(C) =
0$（见\cref{ex:b3:measure:cantor}）：$\mu, \lambda$ 各自
集中于另一\omterm{def:b3:measure:measure}{测度}判为零集的集合上——二者互相
奇异。

\textbf{17.} $\mu(\{x\}) = \lim_{\delta\downarrow0}
\mu(\intoc{x-\delta}x) = \lim(c(x) - c(x-\delta)) = 0$，由
$c$ 的连续性，知无原子。故 $\mu$ 是无原子的
概率\omterm{def:b3:measure:measure}{测度}，集中于一个勒贝格零的紧集上——
既非类似 $\lambda$ 限制的、有密度的连续型，
也非计数\omterm{def:b3:measure:measure}{测度}那样的原子型：这是第三种类型。

\textbf{18.} 片段 $C_\varepsilon$ 对应一个三进区间
$I_\varepsilon$，长为 $3^{-m}$，由第一部分
第 3 题可知，跨越 $I_\varepsilon$ 时，阶梯函数
恰好上升 $2^{-m}$（前 $m$ 位二进制数字取自 $c$
的二进制展开，固定为 $\varepsilon$，其余各位取遍一切可能）。
因此 $\mu(C_\varepsilon) = \mu(I_\varepsilon) =
c(\text{右端点}) - c(\text{左端点}) = 2^{-m}$：深度为 $m$ 的
数字柱\omterm{def:b3:measure:measure}{测度}均为 $2^{-m}$，即 $m$ 枚公平
硬币的分布。

\textbf{19.} 右端定义了博雷尔\omterm{def:b3:measure:measure}{测度}
$\nu = \frac12\,(x \mapsto \tfrac x3)_*\mu + \frac12\,(x
\mapsto \tfrac{x+2}3)_*\mu$ 在 $A$ 上的取值——概率
\omterm{def:b3:measure:measure}{测度}。将 $c$ 整体延拓后，逐点
分情形（$x \leq 0$；三个三分段；$x \geq 1$）验证
恒等式
\[
c(x) = \tfrac12\,c(3x) + \tfrac12\,c(3x - 2),
\]
例如在 $\intcc{1/3}{2/3}$ 上：$\frac12\cdot1 + \frac12\cdot0 =
\frac12 = c(x)$。因此 $\nu$ 在 $\intoc ab$ 上的取值
为 $c(b) - c(a) = \mu(\intoc ab)$，两个有限\omterm{def:b3:measure:measure}{测度}
在半开区间所成的 $\pi$-系上一致，则
它们在 $\mathcal B$ 上重合（见\cref{thm:b3:measure:uniqueness}）：$\nu =
\mu$.

\textbf{20.} 对 $n$ 作归纳：$c_0(1-x) = 1 - c_0(x)$，
若 $c_n$ 具有该对称性，则对 $x \in
\intcc0{1/3}$: $c_{n+1}(1 - x) = \frac12 + \frac12c_n(3(1-x)
- 2) = \frac12 + \frac12c_n(1 - 3x) = \frac12 + \frac12(1 -
c_n(3x)) = 1 - c_{n+1}(x)$；中间三分之一段关于
$\frac12$ 对称；右侧三分之一段即左侧情形的反射。取
极限得 $c(1-x) = 1 - c(x)$。推前\omterm{def:b3:measure:measure}{测度}：$(s_*\mu)(\intoc ab)
= \mu(\intco{1-b}{1-a}) = c(1-a) - c(1-b)$（$\mu$ 无原子，
由第 17 题，端点约定无关紧要）$= (1 -
c(a)) - (1 - c(b)) = \mu(\intoc ab)$：$s_*\mu = \mu$ 由
唯一性。

\textbf{21.} 设 $X \sim \mu$（连续函数关于
$\mu$ 的积分作为黎曼型和的极限而存在，
在深度为 $m$ 的各段上求和（每段\omterm{def:b3:measure:measure}{测度}为 $2^{-m}$），附带
取样误差 $\leq \operatorname{osc} \leq
\norm{f'}_\infty3^{-m}$；这一点将在\cref{ch:b3:lebesgue}
中系统化）。对称性：$1 - X \sim X$，故 $\E X =
\frac12$。自相似性：$X$ 服从 $\frac Y3$ 的分布，概率为
$\frac12$，另以 $\frac{Y+2}3$ 的分布，概率为
$\frac12$，$Y \sim \mu$，于是
\[
\E X^2 = \frac12\,\frac{\E Y^2}9 + \frac12\,
\frac{\E Y^2 + 4\E Y + 4}{9} = \frac{2\E X^2 + 6}{18},
\]
由此得 $\E X^2 = \frac38$，且 $\operatorname{Var}X = \frac38
- \frac14 = \frac18$。均匀分布的
方差为 $\frac1{12} < \frac18$：康托\omterm{def:b3:measure:measure}{测度}集中在端点附近。

\textbf{22.} $C$ 是闭集，且 $\mu(C) = 1$。若\omterm{def:b3:topology:topology}{开集} $I$
与 $C$ 交于 $x$，则深度为 $m$ 的、含 $x$ 的段收缩
到 $x$，故存在 $I_\varepsilon \subseteq I$，且 $\mu(I) \geq
2^{-m} > 0$：任何更小的闭集都不能承载 $\mu$。
$C$ 外的点都有 $\mu$ \omterm{def:b3:measure:measure}{测度}为 $0$ 的间隙\omterm{def:b3:topology:topology}{邻域}，由此
恰得 $\operatorname{supp}\mu = C$。

\textbf{23.} 同一现象，两种表述。第一部分指出，
$c$ 的增长对其导数不可见：$c' = 0$ 几乎处处成立，
所有的上升都集中在零测集 $C$ 上。第五部分
说，\omterm{def:b3:rings:divisibility}{相伴}\omterm{def:b3:measure:measure}{测度} $\mu_c$ 将全部质量都放在同一
零测集上：$\mu_c \perp \lambda$，故不存在密度 $f \geq
0$ 使 $\mu_c(A) = \int_Af\,\dd\lambda$ 成立——密度
必在 $\lambda$ 零测集上为零。对应关系如下：
单调不减有界的 $F \leftrightarrow$ 有限\omterm{def:b3:measure:measure}{测度}
$\mu_F$（第 14--15 题）；$F$ 为其导数的积分
$\leftrightarrow$ $\mu_F$ 有密度（即“绝对
连续”情形）；一般地，$F'$ 几乎处处存在
（对单调 $F$ 成立，见勒贝格微分定理，超出本章范围），
只恢复密度部分。阶梯函数
是极端情形：它连续，几乎处处导数为 $0$ ——
其\omterm{def:b3:measure:measure}{测度}是\emph{纯奇异}的，而微积分基本定理
并非偶然失效，而是恰好差了
$\mu_c(\R) = 1$ 的奇异质量。

\textbf{24.} 固定 $x < y$ 属于 $\intcc01$，并取 $m \geq 0$
使 $3^{-(m+1)} < y - x \leq 3^{-m}$。对任一三进
区间 $\intcc{k3^{-m}}{(k+1)3^{-m}}$ 上，阶梯函数的增量至多为
$2^{-m}$：该区间或为 $2^m$
个 $C_m$ 段之一，此时增量恰为 $2^{-m}$
（第 18 题），否则它含于某阶段的某一间隙的\omterm{def:b3:topology:interior}{闭包}中，
该阶段 $\leq m$，其上 $c$ 为常数。由于 $y - x \leq
3^{-m}$，区间 $\intcc xy$ 至多与两个相邻的
深度为 $m$ 的三进区间相交，故
\[
c(y) - c(x) \leq 2\cdot2^{-m}
= 2\bigl(3^{-m}\bigr)^{s}
< 2\bigl(3(y - x)\bigr)^{s}
= 4\,(y - x)^{s},
\]
利用 $3^{s} = 2$ 与 $3^{-m} < 3(y - x)$。最优性：
$u < v$ 为深度为 $m$ 的康托段的端点，满足 $v - u =
3^{-m}$ 且 $c(v) - c(u) = 2^{-m} = (v - u)^{s}$；若赫尔德
条件 $\abs{c(v) - c(u)} \leq K(v - u)^{t}$（$t > s$）成立时，
必有 $K \geq 2^{-m}3^{mt} = (3^{t}/2)^{m} \to \infty$（因
$3^{t} > 3^{s} = 2$），且 $\intcc01$ 的每个子区间
都包含这样的段，因此失效也是局部的。
\omterm{def:b3:measure:measure}{测度}形式：$\mu\bigl(\intcc{x-r}{x+r}\bigr) \leq
c(x + r) - c(x - r) \leq 4(2r)^{s} = 4\cdot2^{s}r^{s} \leq
8r^{s}$（$c$ 按第 16 题延拓到 $\R$ 上的取值，
$\mu$ 无原子，见第 17 题）。

\textbf{25.} 记 $S_0(x) = \frac x3$ 与 $S_1(x) =
\frac{x+2}3$；假设断言 $\nu = \frac12(S_0)_*\nu +
\frac12(S_1)_*\nu$。迭代 $m$ 次得
\[
\nu = 2^{-m}\sum_{w \in \{0,1\}^m}(S_w)_*\nu,
\qquad S_w = S_{w_1}\circ\dots\circ S_{w_m}.
\]
由于 $\nu$ 承载于 $\intcc01$，且 $S_w(\intcc01) =
I_w$ 即深度为 $m$ 的康托段（由 $w$ 标记），每个
$(S_w)_*\nu$ 是承载于 $I_w$ 的概率\omterm{def:b3:measure:measure}{测度}；
$2^m$ 个段两两不交，是长度为
$3^{-m}$ 的闭区间。固定 $\intoc ab$，令 $N_m$ 为
满足 $I_w \subseteq \intoc ab$ 的段个数。与 $\intoc ab$ 相交
但不含于其中的段必含 $a$ 或 $b$，且一点
至多属于一段，因此
\[
2^{-m}N_m \leq \nu(\intoc ab) \leq 2^{-m}N_m + 2\cdot2^{-m}.
\]
同样的双重不等式对 $\mu$ 也成立（第 19 题给出
相同的迭代），计数同为 $N_m$。于是
$\abs{\nu(\intoc ab) - \mu(\intoc ab)} \leq 2^{-m+1} \to 0$：
$\nu$ 与 $\mu$ 在半开区间的 $\pi$-系上一致，
且二者都是概率\omterm{def:b3:measure:measure}{测度}，故
由\cref{thm:b3:measure:uniqueness}可得 $\nu = \mu$。
阶梯\omterm{def:b3:measure:measure}{测度}因而正是两映射平均格式的\emph{唯一}
不动点——它是 $C$ 的自相似性在\omterm{def:b3:measure:measure}{测度}层面的
表述。
\end{solution}
